\chapter{Requirement Quality Metric Checklists}
\label{appendix:metrics-checklists}

This appendix extends Chapter~\ref{ch:quantifying-quality} with ready-to-use checklists and scorecards. Each table pairs the qualitative guidance from the main text with concrete review prompts, measurement methods, and suggested guardrails so teams can operationalize their quality program.

\section{Completeness and Context Scorecard}
\begin{longtable}{p{0.23\textwidth}p{0.35\textwidth}p{0.3\textwidth}}
\toprule
\textbf{Completeness Element} & \textbf{Review Questions} & \textbf{Measurement / Guardrail}\\
\midrule
Problem Statement & Does the ticket restate the user/business problem in one sentence? & Binary check. Missing problem statement drops score by 15 points.\\
Hypothesis & Are expected outcomes or behavioral changes articulated? & Require explicit ``We expect...'' statement. Missing hypothesis triggers rewrite.\\
Stakeholders & Which personas or systems consume the change? & List all impacted personas; require sign-off from each owner for critical work.\\
Scenarios & Are happy path, failure, and recovery flows described? & Track percentage of tickets covering all three flows; target $\geq80\%$.\\
Data Contracts & Are input/output data tables or schemas referenced? & Capture table or API names. Unreferenced data assets block estimation.\\
Assumptions & Are hidden assumptions and constraints called out? & Require at least one ``Assumptions'' bullet; unanswered item reduces completeness by 10 points.\\
Acceptance Criteria & Is there a numbered list of verifiable outcomes? & Minimum of three acceptance tests or metrics; fewer requires reviewer approval.\\
Risks and Mitigations & Are risks prioritized with mitigation steps? & Require probability/impact tags; unresolved high risk escalated before sprint commit.\\
Metrics / SLO Impact & Which KPIs or \glspl{slo} shift? & Capture baseline and target values. Missing baseline blocks deployment.\\
\bottomrule
\end{longtable}

\section{Clarity and Consistency Checklist}
\begin{longtable}{p{0.26\textwidth}p{0.32\textwidth}p{0.27\textwidth}}
\toprule
\textbf{Dimension} & \textbf{Checklist Prompts} & \textbf{Recommended Tooling}\\
\midrule
Terminology & Are glossary terms (e.g., \gls{slo}, \gls{wape}, \gls{psi}) used consistently? & Glossary lint that flags synonyms or undefined acronyms.\\
Readability & Are sentences $<25$ words and paragraphs $<6$ lines? & Automated readability score with warning below Grade 8 or above Grade 14.\\
Ambiguity & Does the ticket avoid vague qualifiers (``fast,'' ``optimize'')? & Regex/NLP scan producing ``ambiguous term'' count shown in dashboards.\\
Visual Aids & Are diagrams/tables used when prose becomes dense? & Require at least one diagram when readability score exceeds threshold.\\
Dependency Mapping & Are upstream/downstream components explicitly referenced? & Graph view showing each ticket node; gaps trigger review comments.\\
Cross-Doc Consistency & Do linked documents (designs, runbooks) use matching IDs? & Link checker verifying shared IDs, ensuring traceable references.\\
\bottomrule
\end{longtable}

\section{Testability and Predictive Indicators}
\begin{longtable}{p{0.24\textwidth}p{0.35\textwidth}p{0.28\textwidth}}
\toprule
\textbf{Metric} & \textbf{Checklist Items} & \textbf{Threshold / Alert}\\
\midrule
Acceptance-Test Mapping & Does every criterion map to a test case, metric, or alert? & $100\%$ mapping required; flag gaps directly in issue tracker.\\
Instrumentation Plan & Are dashboards, logs, or synthetic tests defined pre-build? & Require named dashboards or log queries before code freeze.\\
Defect Prediction Signal & Does historical data classify the ticket as high risk? & If predicted risk $>0.65$, schedule additional peer review.\\
Rework Probability & Have similar tickets churned mid-sprint due to unclear data? & If rework probability $>0.4$, enforce ``pre-sprint clinic'' session.\\
Cycle-Time Forecast & How does quality score correlate with lead time? & Block sprint entry if forecast indicates delay $>6$ days beyond baseline.\\
Traceability Links & Are OKRs, epics, and test suites linked? & Require at least one upstream OKR and two downstream references.\\
\bottomrule
\end{longtable}

\section{Reviewer Roles and Cadence}
\begin{longtable}{p{0.2\textwidth}p{0.34\textwidth}p{0.34\textwidth}}
\toprule
\textbf{Role} & \textbf{Primary Checklist Focus} & \textbf{Cadence / SLA}\\
\midrule
Product Manager & Completeness, stakeholder coverage, hypothesis clarity. & Submit scorecard before estimation; respond to comments within 1 business day.\\
Engineering Lead & Technical feasibility, dependency graph, instrumentation. & Review within 2 business days of ticket submission.\\
Data Scientist & Metric definitions, \glspl{wape}/\glspl{psi}, data contracts. & Attend weekly quality clinic; log findings in analytics backlog.\\
QA / SDET & Acceptance-test mapping, automation readiness. & Validate checklists before sprint planning; block tickets lacking tests.\\
Program / Portfolio & Traceability, risk register alignment, quality gate compliance. & Run monthly audit sampling 10\% of tickets per product line.\\
\bottomrule
\end{longtable}

Teams can copy these tables into their issue tracker, collaborative docs, or BI dashboards. Update weights, guardrails, and SLAs as quality maturity evolves so the scorecards remain actionable and relevant.
